\documentclass{article}
\usepackage{amsthm,amsmath}

% Definición de entornos
\newtheorem{axiom}{Axioma}
\newtheorem{theorem}{Teorema}
\renewcommand{\proofname}{Prueba}

\begin{document}

\section*{Ejemplo de Axiomas y Teoremas Numerados}

\begin{axiom}\label{ax:1}
Si $a=b$, entonces $a+c=b+c$.
\end{axiom}

\begin{axiom}\label{ax:2}
Si $a=b$ y $b=c$, entonces $a=c$.
\end{axiom}

\begin{axiom}\label{ax:3}
Para todo número real $a$, existe $-a$ tal que $a + (-a) = 0$.
\end{axiom}

\begin{axiom}\label{ax:4}
Existe un número $0$ tal que $a+0=a$ para todo $a$.
\end{axiom}

\begin{theorem}\label{th:1}
Si $a=b$, entonces $a+0=b+0$.
\end{theorem}

\begin{proof}
Por el Axioma~\ref{ax:4}, $a+0=a$ y $b+0=b$.
Aplicando el Axioma~\ref{ax:1}, se obtiene $a+0=b+0$.
\end{proof}

\begin{theorem}\label{th:2}
Si $a=b$, entonces $a+(-a)=b+(-b)=0$.
\end{theorem}

\begin{proof}
Usando los Axiomas~\ref{ax:1}, \ref{ax:3} y \ref{ax:4}, tenemos
\[
a+(-a)=0=b+(-b).
\]
\end{proof}

\end{document}

